% Vietnamese translation for OI Public Library
% Source-PDF: IOI中国国家候选队论文/2009/贾志豪/组合游戏略述——浅谈SG游戏的若干拓展及变形.pdf
\documentclass[11pt]{article}
\usepackage{oipl-vn}

\newcommand{\SG}{\operatorname{SG}}
\newcommand{\mex}{\operatorname{mex}}
\newcommand{\xor}{\mathbin{\oplus}}
\newcommand{\bigxor}{\bigoplus}
\newcommand{\step}{\operatorname{step}}

\title{Lược khảo trò chơi tổ hợp: Một số mở rộng và biến thể của trò chơi SG}
\author{Jia Zhihao\\Trường Trung học số 2 Thạch Gia Trang}
\date{2009}

\begin{document}
\maketitle

\origtitle{Lược khảo trò chơi tổ hợp: Một số mở rộng và biến thể của trò chơi SG}
\sourcepath{IOI China candidate team papers / 2009 / Jia Zhihao}

\begin{abstract}
Bài viết bàn về một lớp trò chơi tổ hợp có thể xử lý bằng hàm Sprague--Grundy
(SG), rồi trình bày nhiều biến thể thường gặp: trò chơi anti-SG, multi-SG,
every-SG, trò lật đồng xu và trò xóa cạnh trên đồ thị. Trọng tâm không phải
là kể lại các bài toán có bối cảnh dài, mà là tách ra mô hình toán học và
chứng minh vì sao các mô hình đó quy về Nim hoặc biến thể của Nim.
\end{abstract}

\paragraph{Từ khóa.}
Hàm SG; mô hình Nim; trò chơi anti; trò chơi multi; trò chơi every; lật đồng
xu; xóa cạnh; chứng minh; mở rộng.

\tableofcontents

\section{Nền tảng}

\subsection{Phạm vi thảo luận}

Bài viết xét các trò chơi tổ hợp thỏa những điều kiện sau:
\begin{itemize}
  \item Có hai người chơi, luân phiên quyết định, và cả hai đều chơi tối ưu.
  \item Khi một người không thể đi, trò chơi kết thúc và người đó thua. Dù
  hai bên đi thế nào, trò chơi cũng kết thúc sau hữu hạn bước.
  \item Một trạng thái không thể được đạt tới nhiều lần; không có hòa.
  \item Tập nước đi hợp lệ tại một trạng thái chỉ phụ thuộc trạng thái đó,
  không phụ thuộc người chơi.
\end{itemize}
Những trò chơi như vậy được gọi ở đây là \term{trò chơi tổ hợp SG}.

``Tối ưu'' nghĩa là: nếu người đang đi có một nước đưa đối thủ vào trạng thái
chắc chắn thua, người đó sẽ chọn một nước như vậy.

\subsection{Chuyển không gian trạng thái thành đồ thị}

Mỗi trạng thái trò chơi có thể xem như một đỉnh của đồ thị, mỗi quyết định là
một cạnh có hướng. Đồ thị này gọi là \term{đồ thị trò chơi}. Một ván chơi khi
đó là một đường đi từ một đỉnh ban đầu tới một đỉnh có bậc ra bằng \(0\).
Cách nhìn này giúp bỏ bối cảnh cụ thể và nhấn mạnh bản chất toán học.

\subsection{Hàm SG}

Hàm SG là hàm đánh giá từng đỉnh trong đồ thị trò chơi:
\[
  \SG(v)=\mex\{\SG(u)\mid \text{có cạnh từ } v \text{ tới } u\}.
\]
Ở đây \(\mex(A)\), minimal excludant, là số tự nhiên nhỏ nhất không thuộc
tập \(A\):
\[
  \mex(A)=\min\{k\mid k\notin A,\ k\in\mathbb{N}\}.
\]

Về mặt lý thuyết, ta có thể tính mọi giá trị SG theo thứ tự topo của đồ thị,
rồi chỉ cần biết \(\SG(v)\) có bằng \(0\) hay không để xác định thắng thua.
Tuy nhiên cách này thường quá tốn thời gian và bộ nhớ vì chưa tận dụng cấu
trúc đặc biệt của trò chơi.

Hai tính chất cơ bản của SG là:
\begin{enumerate}
  \item Nếu \(\SG(v)=0\), mọi hậu duệ trực tiếp của \(v\) đều có SG khác \(0\).
  \item Nếu \(\SG(v)\ne 0\), tồn tại một hậu duệ trực tiếp có SG bằng \(0\).
\end{enumerate}

\subsection{Mô hình Nim}

\term{Tổng của các trò chơi} là trò chơi tạo bởi nhiều trò SG độc lập: mỗi
lượt người chơi chọn đúng một trò con chưa kết thúc và đi một nước trong trò
đó; người không thể đi thua.

Trong Nim, có \(N\) đống đá, mỗi lượt chọn một đống và lấy đi một số đá dương;
người lấy viên cuối thắng. Với một đống có \(k\) viên, giá trị SG là \(k\).
Với nhiều đống, tổng SG là xor:
\[
  \SG_{\mathrm{tot}}=\SG_1\xor\SG_2\xor\cdots\xor\SG_n.
\]
Người đi trước thắng khi và chỉ khi giá trị này khác \(0\).

\paragraph{Định lý.}
Trong tổng của các trò chơi SG, nếu thay bất kỳ trò con nào bằng một đống Nim
có kích thước đúng bằng giá trị SG của trò đó, kết quả thắng thua không đổi.
Nói cách khác, khi xét tổng trò chơi, mọi chi tiết của trò con có thể được
nén vào một số SG.

\section{Các biến thể luật}

\subsection{Người đi nước cuối thua: Anti-SG và định lý SJ}

Xét trước anti-Nim:
\begin{itemize}
  \item Có \(N\) đống đá.
  \item Mỗi lượt lấy một số đá dương từ đúng một đống.
  \item Người lấy viên cuối cùng thua.
\end{itemize}

Trong anti-Nim, người đi trước thắng khi và chỉ khi:
\begin{enumerate}
  \item mọi đống đều có kích thước \(1\) và xor của các đống bằng \(0\); hoặc
  \item tồn tại đống có kích thước lớn hơn \(1\) và xor của các đống khác \(0\).
\end{enumerate}
Nếu mọi đống đều là \(1\), thắng thua chỉ phụ thuộc vào chẵn lẻ số đống. Nếu
có đống lớn hơn \(1\), khi xor khác \(0\) người đi trước hoặc đưa xor về \(0\),
hoặc nếu chỉ còn một đống lớn thì biến trạng thái thành số lẻ đống \(1\).
Khi xor bằng \(0\) và còn ít nhất hai đống lớn, mọi nước đi đều đưa đối thủ
vào trạng thái thắng.

Lập luận anti-Nim không thể áp dụng máy móc cho mọi trò SG, vì trong đồ thị
SG tổng quát, trạng thái SG bằng \(0\) không nhất thiết là trạng thái kết thúc.
Do đó bài viết đưa ra một phiên bản có điều kiện, gọi là định lý SJ.

\paragraph{Định nghĩa.}
Trong \term{Anti-SG}, người chơi có tập nước đi rỗng là người thắng; các quy
tắc khác giống trò SG thường.

\paragraph{Định lý SJ.}
Xét một Anti-SG và quy định thêm rằng khi mọi trò con đều có SG bằng \(0\) thì
trò chơi kết thúc. Khi đó người đi trước thắng khi và chỉ khi:
\begin{enumerate}
  \item xor SG tổng khác \(0\) và tồn tại một trò con có SG lớn hơn \(1\); hoặc
  \item xor SG tổng bằng \(0\) và không có trò con nào có SG lớn hơn \(1\).
\end{enumerate}

\paragraph{Ý chứng minh.}
Cần kiểm tra ba điều: mọi trạng thái kết thúc là thắng cho người sắp đi; mọi
trạng thái thua chỉ đi được tới trạng thái thắng; và mọi trạng thái thắng có
một nước đi tới trạng thái thua.

Nếu xor bằng \(0\) và có trò con SG lớn hơn \(1\), thực ra phải có ít nhất hai
trò con như vậy. Mỗi nước chỉ đổi một trò con, nên sau một nước vẫn còn một
trò con SG lớn hơn \(1\), đồng thời xor mới khác \(0\); đó là trạng thái thắng
cho đối thủ.

Nếu xor khác \(0\) nhưng mọi trò con chỉ có SG \(0\) hoặc \(1\), thì số trò
con có SG \(1\) là lẻ. Mọi nước đi hoặc tạo đúng một trò con SG lớn hơn \(1\)
với xor khác \(0\), hoặc biến chẵn/lẻ của các trò SG \(1\) để tạo trạng thái
xor bằng \(0\) không có SG lớn hơn \(1\). Cả hai đều là trạng thái thắng cho
đối thủ.

Hai trường hợp thắng còn lại có nước đi chuẩn: nếu tồn tại trò con SG lớn hơn
\(1\), dùng tính chất SG để đưa xor về \(0\) mà vẫn còn trò con lớn hơn \(1\);
nếu mọi trò con chỉ có SG \(0/1\) và xor bằng \(0\), nhưng chưa kết thúc, đổi
một trò con SG \(1\) thành \(0\), tạo số lẻ trò con SG \(1\).

\subsection{Một đống có thể tách thành nhiều đống: Multi-SG}

\paragraph{Định nghĩa.}
Trong \term{Multi-SG}, dưới điều kiện vẫn tôn trọng thứ tự topo của đồ thị
trò chơi, hậu duệ của một trò con đơn lẻ có thể là nhiều trò con mới. Các quy
tắc khác giống SG thường.

Vì đồ thị trạng thái vẫn có thứ tự topo, ta vẫn tính SG bằng quy nạp theo thứ
tự đó. Khi một nước biến một trò thành nhiều trò, giá trị của hậu duệ chính
là xor các giá trị SG của những trò mới. Do đó khung SG vẫn hoạt động; điểm
khác biệt chỉ nằm ở cách liệt kê hậu duệ.

\subsection{Mọi quân có thể đi đều phải đi: Every-SG}

\paragraph{Định nghĩa.}
Trong \term{Every-SG}, với mọi trò con chưa kết thúc, người chơi bắt buộc phải
đi một nước trong trò con đó. Các quy tắc khác giống trò SG thường.

Ở đây không chỉ có thắng/thua, mà còn có dài/ngắn. Người đang ở thế thắng muốn
kéo dài các trò con thắng của mình đủ lâu, trong khi đối thủ muốn kết thúc
chúng sớm.

Khi tính một trạng thái \(v\), ngoài việc biết \(\SG(v)\) có bằng \(0\) hay
không, ta tính thêm
\[
\step(v)=
\begin{cases}
0, & v \text{ là trạng thái kết thúc},\\
\max\{\step(u)\}+1, & \SG(v)>0,\ u \text{ là hậu duệ và } \SG(u)=0,\\
\min\{\step(u)\}+1, & \SG(v)=0,\ u \text{ là hậu duệ}.
\end{cases}
\]

\paragraph{Định lý.}
Trong Every-SG, người đi trước thắng khi và chỉ khi giá trị \(\step\) lớn
nhất trong các trò con là số lẻ.

\paragraph{Ý chứng minh.}
Bằng quy nạp trên đồ thị, trạng thái thắng có \(\step\) lẻ, trạng thái thua
có \(\step\) chẵn. Lấy trò con \(A\) có \(\step\) lớn nhất. Người thắng luôn
đi trong \(A\) khi \(A\) đang là trạng thái thắng và có thể đảm bảo mỗi lượt
chỉ làm \(\step(A)\) giảm tối đa \(1\). Đối thủ khi chạm vào \(A\) cũng chỉ
có thể làm nó giảm hữu hạn theo quy tắc tương ứng. Với mọi trò con khác mà
người thắng đang ở thế thua, người thắng có thể buộc chúng kết thúc không
muộn hơn \(A\). Vì vậy chẵn/lẻ của \(\step_{\max}\) quyết định thắng thua.

\section{Một số mô hình thường gặp}

\subsection{Trò lật đồng xu}

Luật tổng quát:
\begin{itemize}
  \item Có \(N\) đồng xu xếp thành hàng, mỗi đồng có thể ngửa hoặc sấp.
  \item Người chơi lật các đồng xu theo một số ràng buộc, ví dụ mỗi lần lật
  một hoặc hai đồng, hoặc lật một đoạn liên tiếp.
  \item Trong các đồng bị lật, đồng bên phải nhất bắt buộc phải đổi từ ngửa
  sang sấp.
  \item Ai không thể lật thì thua.
\end{itemize}

Kết luận quan trọng: giá trị SG của toàn cục bằng xor các giá trị SG của từng
đồng đang ngửa nếu nó tồn tại một mình:
\[
  \SG(\text{trạng thái})=\bigxor_{\text{đồng }i\text{ ngửa}}\SG(i).
\]
Ý chứng minh là quy nạp theo trọng số trạng thái. Gán đồng thứ \(k\) từ trái
sang trọng số \(2^k\), và trọng số trạng thái là tổng trọng số các đồng ngửa.
Mọi nước đi làm giảm trọng số này vì đồng ngửa bên phải nhất bị lật xuống.
Một nước đi có thể xem như bỏ một ``đống'' tương ứng với đồng bên phải nhất,
rồi thêm vào một số hậu duệ ở bên trái; điều này tương đương với nước đi trong
Nim.

\subsection{Trò xóa cạnh trên cây}

Luật:
\begin{itemize}
  \item Cho một cây có gốc.
  \item Mỗi lượt xóa một cạnh; phần không còn nối với gốc bị xóa theo.
  \item Ai không có nước đi thì thua.
\end{itemize}

\paragraph{Định lý.}
Lá có SG bằng \(0\). Với một đỉnh trong, giá trị SG của nó là xor của
\(\SG(\text{con})+1\) trên mọi con:
\[
  \SG(v)=\bigxor_{u\text{ là con của }v}(\SG(u)+1).
\]

\paragraph{Ý chứng minh.}
Nếu gốc \(A\) chỉ có một con \(B\), gọi \(G'\) là cây con gốc \(B\). Xóa cạnh
\(AB\) đưa trò chơi về \(0\); xóa một cạnh bên trong \(G'\) tương ứng với việc
đưa \(G'\) tới một hậu duệ rồi cộng thêm cạnh \(AB\). Do đó
\[
  \SG(G)=\SG(G')+1.
\]
Nếu gốc có nhiều con, các cây con là các trò độc lập trong tổng trò chơi, nên
giá trị toàn cục là xor các giá trị của từng nhánh.

\subsection{Christmas Game}

Trong bài PKU3710, có nhiều đồ thị liên thông cục bộ. Hai người luân phiên
xóa cạnh; phần tách khỏi gốc bị loại bỏ. Đồ thị được tạo từ cây cơ sở bằng
cách thêm một số cạnh, các chu trình không dùng chung cạnh và mỗi chu trình
chỉ có một điểm chung với cây cơ sở.

Điểm khác với cây là có chu trình. Vì mọi chu trình đều là một vòng đơn gắn
vào một đỉnh của cây, ta phân tích theo độ dài chu trình:
\begin{itemize}
  \item Chu trình lẻ có SG bằng \(1\): xóa một cạnh sẽ tạo hai dây có cùng
  chẵn/lẻ, xor của chúng không thể nhận giá trị lẻ tương ứng.
  \item Chu trình chẵn có SG bằng \(0\): xóa một cạnh sẽ tạo hai dây khác
  chẵn/lẻ, xor không thể bằng \(0\).
\end{itemize}
Vì vậy có thể bỏ chu trình chẵn và thay chu trình lẻ bằng một dây độ dài \(1\).
Sau biến đổi, bài toán trở lại trò xóa cạnh trên cây; tính SG từng cây rồi xor
theo mô hình Nim.

\subsection{Trò xóa cạnh trên đồ thị vô hướng}

Bỏ điều kiện đặc biệt về chu trình, ta có trò chơi trên đồ thị vô hướng liên
thông có gốc: mỗi lượt xóa một cạnh, phần không nối với gốc bị xóa, ai không
đi được thì thua. Công cụ chuẩn là \term{Fusion Principle}.

\paragraph{Fusion Principle.}
Trong đồ thị vô hướng, có thể co một chu trình chẵn thành một đỉnh mới; một
chu trình lẻ thành một đỉnh mới cộng thêm một cạnh mới; mọi cạnh nối vào chu
trình cũ được nối vào đỉnh mới. Biến đổi này không đổi giá trị SG.

Chứng minh nguyên lý này khá dài và được trình bày trong \textit{Winning Ways}.
Sau khi áp dụng nó lặp lại, đồ thị vô hướng có thể được đưa về cấu trúc cây,
rồi dùng công thức của trò xóa cạnh trên cây.

\section{Tài liệu tham khảo}

\begin{enumerate}
  \item Thomas S. Ferguson, \textit{Game Theory}.
  \item Zhang Yifei, luận văn đội tuyển quốc gia 2002,
  \textit{Từ nhận thức cảm tính đến nhận thức lý tính: phân tích lời giải một
  lớp trò chơi đối kháng}.
  \item Wang Xiaoke, luận văn đội tuyển quốc gia 2007,
  \textit{Phân tích một lớp trò chơi tổ hợp}.
  \item Luo Ji, luận văn đội tuyển quốc gia 2002,
  \textit{Hai hướng suy nghĩ để giải bài toán đối sách}.
  \item E. R. Berlekamp, J. H. Conway, R. K. Guy,
  \textit{Winning Ways for Your Mathematical Plays}, tập 1.
  \item Tài liệu \textit{Game Strategy} của thầy Zhu Quanmin.
\end{enumerate}

\section{Phụ lục: bài luyện tập}

\begin{itemize}
  \item \textit{Strips}, POI 2000.
  \item \textit{Christmas Game}, PKU 3710.
  \item \textit{Georgia and Bob}, PKU 1704.
  \item \textit{Crosses and Crosses}, PKU 3537.
  \item \textit{KNIGHT}, CEOI 2008.
  \item \textit{NIM}, PKU 2608.
  \item \textit{NIM}, PKU 2975.
  \item \textit{S-NIM}, PKU 2960.
  \item \textit{2D-NIM}, PKU 1021.
  \item \textit{I Love this Game}, PKU 1678.
\end{itemize}

\end{document}
